\section{Transformer-Based Neural Operator}

\subsection{Motivation for an attention-based surrogate}

The coupled thermoelastic system written down in Chapter~2 defines a
deterministic mapping
\begin{equation}
\mathcal{G}: (a,\,\mathrm{IC},\,\mathrm{BC},\,s) \;\longmapsto\; \big(T(x,t),\,u(x,t),\,\sigma(x,t)\big),
\end{equation}
where the input collects the material parameter field
\(a(x) = (\rho,\lambda,\mu,\gamma,k,C_p)\), the initial and
boundary conditions, and the source description. A neural surrogate
\(\mathcal{G}_\theta\) with learnable parameters \(\theta\) is intended
to approximate this mapping at a fraction of the cost of running the
COMSOL solver of Chapter~4. Three properties of the present problem
constrain the choice of architecture. First, the geological domain is
discretised by an unstructured tetrahedral mesh and not by a regular
grid, so the surrogate must accept arbitrary point clouds. Second, the
thermoelastic response is non-local: a thermal pulse in one region of
the rock eventually influences distant nodes through both heat
diffusion and elastic wave propagation, and the surrogate has to relate
sample points that may be far apart in physical space. Third, training
and inference should ideally operate on \emph{different} samplings of
the domain, since the available number of nodes is dictated by the
COMSOL mesh whereas a user of the predictor may ask for a value at any
spatial coordinate. The Transformer architecture, originally introduced
for sequence modelling \cite{vaswani2017}, naturally satisfies all
three constraints when it is reformulated as a neural operator acting
on a set of points \cite{liNeuralOperator2023}. It treats the
discretisation as a finite token set rather than as a regular array,
its attention mechanism couples every pair of tokens directly, and the
output is queried at locations that need not coincide with the input
sampling. For these reasons we adopt a Transformer-based neural
operator as one of the four surrogate families compared in this
thesis.

\subsubsection{Tokenisation of the thermoelastic state}

The first step is to encode the physical state of the medium as a
sequence of \emph{tokens}. For every retained mesh node
\(x_i\in\Omega\) we form an input token whose feature vector
combines the local geometry, the initial physical state, and the
material parameters of the medium:
\begin{equation}
a_i =
\big[\,
  x_i,\,y_i,\;
  T_i^{(0)},\;
  u_i^{(0)},\,v_i^{(0)},\;
  \dot u_i^{(0)},\,\dot v_i^{(0)},\;
  E_i,\,\nu_i,\,\rho_i,\,\alpha_i,\,k_i,\,C_{p,i}
\,\big]\in\mathbb{R}^{13}.
\end{equation}
The first two entries fix the location at which the token is
attached, the next five entries describe the state of the rock at the
beginning of the simulation interval, and the remaining six entries are
the static material constants of the medium. A separate set of
\emph{query tokens} \(\{q_j\}_{j=1}^{N_q}\) is then formed from the
locations at which the predicted fields are desired:
\begin{equation}
q_j = (x_q^{(j)},\,y_q^{(j)})\in\mathbb{R}^{2}.
\end{equation}
The number of input tokens \(N_{\mathrm{in}}\) and the number of query
tokens \(N_q\) are independent of each other and may differ between
training and inference; this is the property that makes the
construction discretisation-invariant.

\begin{figure}[ht]
\centering
\begin{tikzpicture}[
  font=\small,
  node distance=4mm and 14mm,
  itok/.style={
    rectangle, rounded corners=2pt,
    draw=blue!55!black, fill=blue!8,
    minimum width=64mm, minimum height=8mm,
    inner sep=3pt, align=left
  },
  qtok/.style={
    rectangle, rounded corners=2pt,
    draw=teal!55!black, fill=teal!10,
    minimum width=42mm, minimum height=8mm,
    inner sep=3pt, align=left
  }
]

\node[itok] (i1) {\(a_1 = (x_1,\;T_1^{(0)},\;u_1^{(0)},\;\dot{u}_1^{(0)},\;a_1^{\text{mat}})\)};
\node[itok, below=of i1] (i2) {\(a_2 = (x_2,\;\ldots)\)};
\node[below=2mm of i2] (idots) {\(\vdots\)};
\node[itok, below=2mm of idots] (iN) {\(a_{N_{\mathrm{in}}} = (x_{N_{\mathrm{in}}},\;\ldots)\)};

\node[draw=black!50, dashed, fit=(i1) (iN), inner sep=4mm, label=above:{Input set (mesh nodes, ground-truth state at \(t=0\))}] (inset) {};

\node[qtok, right=18mm of i1.east] (q1) {\(q_1 = (x_q^{(1)})\)};
\node[qtok, below=of q1] (q2) {\(q_2 = (x_q^{(2)})\)};
\node[below=2mm of q2] (qdots) {\(\vdots\)};
\node[qtok, below=2mm of qdots] (qN) {\(q_{N_q} = (x_q^{(N_q)})\)};

\node[draw=black!50, dashed, fit=(q1) (qN), inner sep=4mm, label=above:{Query set (arbitrary points)}] (qset) {};

\end{tikzpicture}
\caption{Tokenisation of the thermoelastic state used by the
Transformer-based neural operator. Input tokens carry their spatial
coordinates together with the initial thermal and mechanical state of
the rock and the local material parameters. Query tokens carry only
coordinates and indicate where the predicted fields should be returned.}
\label{fig:transformer-tokens}
\end{figure}

\subsection{Scaled dot-product attention}

The basic operation in a Transformer is attention. Given a query
representation \(Q\in\mathbb{R}^{N_q\times d_k}\), key and
value representations \(K\in\mathbb{R}^{N_{\mathrm{in}}\times d_k}\)
and \(V\in\mathbb{R}^{N_{\mathrm{in}}\times d_v}\), the
scaled dot-product attention is defined as
\begin{equation}
\operatorname{Attention}(Q,K,V) =
\operatorname{softmax}\!\left(
  \frac{QK^{\top}}{\sqrt{d_k}}
\right)V.
\label{eq:dotproduct}
\end{equation}
The matrix \(QK^{\top}\in\mathbb{R}^{N_q\times N_{\mathrm{in}}}\)
contains a similarity score between every query and every input. The
scaling factor \(\sqrt{d_k}\) prevents the dot products from growing
with the embedding dimension and keeps the softmax in its informative
regime. The row-wise softmax converts the scores into a non-negative
weight matrix whose rows sum to one. Multiplying these weights by
\(V\) returns, for each query, a convex combination of input
values that is selected by the learnt similarity pattern. From the
operator-learning point of view this operation is a soft, data-driven
analogue of an integral kernel, in which the kernel itself depends on
the input through its keys and queries.

\subsection{Multi-head attention}

A single attention head describes one relational pattern. To express
several patterns in parallel, the embedding space is split into
\(h\) independent subspaces. For each head \(i\in\{1,\dots,h\}\) the
tokens are projected by learnt matrices and a separate attention is
computed:
\begin{equation}
\operatorname{head}_i =
\operatorname{Attention}\!\big(
  QW^Q_i,\,
  KW^K_i,\,
  VW^V_i
\big),
\end{equation}
\begin{equation}
\operatorname{MultiHead}(Q,K,V) =
\operatorname{Concat}(\operatorname{head}_1,\dots,\operatorname{head}_h)\,W^O,
\end{equation}
where the projection matrices satisfy
\(W^Q_i,W^K_i,W^V_i\in\mathbb{R}^{d_m\times d_k}\)
and \(W^O\in\mathbb{R}^{h d_v\times d_m}\). In the baseline
configuration used in this thesis the embedding dimension is
\(d_m=128\), the number of heads is \(h=4\), and the per-head key and
value dimensions are \(d_k=d_v=d_m/h=32\). Different heads can in
principle specialise in different physical interactions, for example
near-field heat diffusion in one head and longer-range elastic
coupling in another.

\subsection{Encoder block}

The input tokens \(\{a_i\}\) are first lifted into the embedding space
by a learnt linear projection
\(W_{\mathrm{in}}\in\mathbb{R}^{C_{\mathrm{in}}\times d_m}\),
yielding initial representations
\(H^{(0)}\in\mathbb{R}^{N_{\mathrm{in}}\times d_m}\) with
\(C_{\mathrm{in}}=13\). They are then refined by a stack of
\(L_{\mathrm{enc}}\) identical encoder blocks. Each block consists of
multi-head self-attention followed by a position-wise feed-forward
network, with residual connections and layer normalisation around
both:
\begin{equation}
H^{(\ell)'} =
\operatorname{LayerNorm}\!\big(
  H^{(\ell)} +
  \operatorname{MultiHead}(H^{(\ell)},H^{(\ell)},H^{(\ell)})
\big),
\end{equation}
\begin{equation}
H^{(\ell+1)} =
\operatorname{LayerNorm}\!\big(
  H^{(\ell)'} +
  \operatorname{FFN}(H^{(\ell)'})
\big).
\end{equation}
The feed-forward sub-block applies the same two-layer perceptron to
every token independently,
\begin{equation}
\operatorname{FFN}(h) =
W_2\,\phi\!\big(W_1h + b_1\big) + b_2,
\end{equation}
where \(W_1\in\mathbb{R}^{d_{\mathrm{ff}}\times d_m}\),
\(W_2\in\mathbb{R}^{d_m\times d_{\mathrm{ff}}}\), and the
activation \(\phi\) is taken to be the GELU function. The hidden width
of the FFN sub-block is \(d_{\mathrm{ff}} = 4\,d_m\), in agreement with
the original Transformer choice \cite{vaswani2017}. After
\(L_{\mathrm{enc}}\) such layers the encoder produces a context tensor
\(H\in\mathbb{R}^{N_{\mathrm{in}}\times d_m}\) that captures
how each input token relates to every other input token through the
learnt self-attention pattern.

\subsection{Cross-attention decoder for query points}

The decoder is the operator-learning component. Its purpose is to
compute the predicted state at every query location by attending to
the encoded inputs. The query coordinates are first lifted by a
separate linear projection
\(W_{\mathrm{q}}\in\mathbb{R}^{2\times d_m}\) into a query
embedding \(Z^{(0)}\in\mathbb{R}^{N_q\times d_m}\). A stack of
\(L_{\mathrm{dec}}\) decoder blocks then refines this representation
via cross-attention between the query embedding (as queries) and the
encoder output (as keys and values):
\begin{equation}
Z^{(\ell)'} =
\operatorname{LayerNorm}\!\big(
  Z^{(\ell)} +
  \operatorname{MultiHead}(Z^{(\ell)},H,H)
\big),
\end{equation}
\begin{equation}
Z^{(\ell+1)} =
\operatorname{LayerNorm}\!\big(
  Z^{(\ell)'} +
  \operatorname{FFN}(Z^{(\ell)'})
\big).
\end{equation}
After the final decoder layer a linear output head
\(W_{\mathrm{out}}\in\mathbb{R}^{d_m\times C_{\mathrm{out}}}\)
maps every decoded query token to a vector of predicted physical
quantities:
\begin{equation}
\hat{y}_j =
\big(\hat T_j,\,\hat u_j,\,\hat v_j\big) =
W_{\mathrm{out}}\,Z^{(L_{\mathrm{dec}})}_{j,:}
\in\mathbb{R}^{C_{\mathrm{out}}},
\qquad C_{\mathrm{out}} = 3.
\end{equation}
Within the framework of Chapter~2 these three channels coincide with
the primary unknowns of the two-dimensional thermoelastic system: the
temperature field and the two in-plane components of the displacement
vector. The other
quantities of physical interest, in particular the strain and stress
tensors, can either be derived from the predicted displacement by
spatial differentiation or be added as additional output channels in
future iterations of the model.

\subsection{Architecture used in this work}

Figure~\ref{fig:transformer-arch} shows the architecture in its
complete form. The hyperparameters retained for the baseline of this
thesis are \(d_m=128\), \(h=4\),
\(L_{\mathrm{enc}}=L_{\mathrm{dec}}=4\), \(d_{\mathrm{ff}}=4\,d_m\),
dropout~\(0.1\), and GELU activation in the feed-forward sub-blocks.
The resulting model contains on the order of \(2.4\times 10^6\)
trainable parameters, which is small by current standards and was
chosen so that training and inference remain feasible on a CPU.

\begin{figure}[ht]
\centering
\begin{tikzpicture}[
  font=\small,
  node distance=3mm and 12mm,
  inp/.style={
    rectangle, rounded corners=2pt,
    draw=blue!55!black, fill=blue!8,
    minimum width=48mm, minimum height=9mm, align=center
  },
  enc/.style={
    rectangle, rounded corners=2pt,
    draw=violet!55!black, fill=violet!10,
    minimum width=48mm, minimum height=9mm, align=center
  },
  qry/.style={
    rectangle, rounded corners=2pt,
    draw=teal!55!black, fill=teal!10,
    minimum width=42mm, minimum height=9mm, align=center
  },
  dec/.style={
    rectangle, rounded corners=2pt,
    draw=orange!70!black, fill=orange!10,
    minimum width=42mm, minimum height=9mm, align=center
  },
  outp/.style={
    rectangle, rounded corners=2pt,
    draw=green!50!black, fill=green!12,
    minimum width=42mm, minimum height=9mm, align=center
  },
  arr/.style={-{Latex[length=2mm]}, thick, draw=black!65}
]

\node[inp] (in)   {Input tokens \(\{a_i\}\in\mathbb{R}^{13}\)};
\node[inp, below=of in] (ip) {Input projection \(\mathbb{R}^{13}\!\to\!\mathbb{R}^{d_m}\)};
\node[enc, below=of ip] (e1) {Self-attention layer \#1};
\node[enc, below=of e1] (e2) {Self-attention layer \#2};
\node[below=2mm of e2] (edots) {\(\vdots\)};
\node[enc, below=2mm of edots] (eL) {Self-attention layer \(L_{\mathrm{enc}}\)};

\node[qry, right=24mm of in] (q)  {Query tokens \(\{q_j\}\in\mathbb{R}^{2}\)};
\node[qry, below=of q]   (qp) {Query projection \(\mathbb{R}^{2}\!\to\!\mathbb{R}^{d_m}\)};
\node[dec, below=of qp]  (d1) {Cross-attention layer \#1};
\node[dec, below=of d1]  (d2) {Cross-attention layer \#2};
\node[below=2mm of d2]   (ddots) {\(\vdots\)};
\node[dec, below=2mm of ddots] (dL) {Cross-attention layer \(L_{\mathrm{dec}}\)};

\node[outp, below=18mm of $(eL)!0.5!(dL)$] (head) {Linear head \(\mathbb{R}^{d_m}\!\to\!\mathbb{R}^{2}\)};
\node[outp, below=of head] (yhat) {\(\hat{y}_j = (\hat T_j,\,\hat u_j,\,\hat v_j)\)};

\draw[arr] (in) -- (ip);
\draw[arr] (ip) -- (e1);
\draw[arr] (e1) -- (e2);
\draw[arr] (e2) -- (edots);
\draw[arr] (edots) -- (eL);

\draw[arr] (q)  -- (qp);
\draw[arr] (qp) -- (d1);
\draw[arr] (d1) -- (d2);
\draw[arr] (d2) -- (ddots);
\draw[arr] (ddots) -- (dL);

\draw[arr] (eL.east) -- ++(8mm,0) |- (d1.west) node[pos=0.25, right] {\scriptsize \(K,\,V\)};
\draw[arr, dashed, draw=black!40] (eL.east) -- ++(8mm,0) |- (dL.west);

\draw[arr] (dL.south) -- ++(0,-4mm) -| (head.north);
\draw[arr] (head) -- (yhat);

\end{tikzpicture}
\caption{End-to-end architecture of the Transformer-based neural
operator. The left column processes the input set with
\(L_{\mathrm{enc}}\) self-attention layers. The right column lifts the
query coordinates and applies \(L_{\mathrm{dec}}\) cross-attention
layers in which the encoded context provides keys and values. A linear
head maps the decoded query embeddings to the three predicted physical
channels at every query location.}
\label{fig:transformer-arch}
\end{figure}

\subsection{Training objective}

The Transformer is trained on COMSOL ground-truth solutions of the
coupled thermoelastic problem of Chapter~4. The primary supervised
loss is the mean squared error between the predicted and the reference
fields, evaluated on the set of query tokens that are randomly drawn
at every training step:
\begin{equation}
\mathcal{L}_{\mathrm{sup}}(\theta) =
\frac{1}{N_q\,C_{\mathrm{out}}}
\sum_{j=1}^{N_q}\sum_{c=1}^{C_{\mathrm{out}}}
\big(\hat{y}_{j,c} - y_{j,c}\big)^2.
\label{eq:sup}
\end{equation}
A complementary \emph{velocity-consistency} term aligns the time
derivative of the predicted displacement with the COMSOL velocity
field. Since the displacement and its time derivative are available as
two distinct exports from the solver, this term acts as an auxiliary
regulariser that anchors the rate of change of the model output and
not just its instantaneous value:
\begin{equation}
\mathcal{L}_{\mathrm{vel}}(\theta) =
\frac{1}{N_q}\sum_{j=1}^{N_q}
\big\|\,\partial_t \hat{u}_j - \dot{u}_j^{\mathrm{ref}}\big\|_2^2,
\label{eq:vel}
\end{equation}
where the time derivative \(\partial_t\hat{u}_j\) is obtained
by automatic differentiation through the model output with respect to
the query time coordinate. A third term enforces the diffusion side of
the heat equation derived in Chapter~2,
\begin{equation}
\mathcal{L}_{\mathrm{th}}(\theta) =
\frac{1}{N_q}\sum_{j=1}^{N_q}
\bigg(
  \partial_t \hat T_j
  - \mathfrak{a}_j\,\nabla^2 \hat T_j
\bigg)^{\!2},
\qquad
\mathfrak{a}_j = \frac{k_j}{\rho_j C_{p,j}},
\label{eq:th}
\end{equation}
which couples the temperature prediction to the local thermal
diffusivity. Equations~\eqref{eq:sup}--\eqref{eq:th} are combined into
a single training objective with non-negative weights:
\begin{equation}
\mathcal{L}(\theta) =
\lambda_{\mathrm{sup}}\mathcal{L}_{\mathrm{sup}}
+ \lambda_{\mathrm{vel}}\mathcal{L}_{\mathrm{vel}}
+ \lambda_{\mathrm{th}}\mathcal{L}_{\mathrm{th}}.
\end{equation}
The default weights used in our experiments are
\((\lambda_{\mathrm{sup}},\,\lambda_{\mathrm{vel}},\,\lambda_{\mathrm{th}}) = (1.0,\;0.25,\;0.05)\).
The objective is minimised with the AdamW optimiser, a cosine
learning-rate schedule, gradient clipping at norm one, and a fixed
random seed for reproducibility. To keep both the training cost and
the memory footprint quadratic in a small token count rather than in
the full mesh size, we sample \(N_{\mathrm{in}} = N_q = 1024\) tokens
uniformly at random from the COMSOL nodes at every training step. This
\emph{token subsampling} simultaneously plays two roles: it controls
the cost of attention and it enforces discretisation invariance, since
the model never sees the full mesh during a single forward pass.

\subsection{Comparison with related Transformer neural operators}

The architecture above is intentionally the simplest member of the
family of Transformer neural operators. Three more elaborate variants
are commonly cited in the recent operator-learning literature and are
summarised in Table~\ref{tab:transformer-variants}. The OFormer
\cite{liNeuralOperator2023} formalises the use of separate input and
query token sets and demonstrates that the cross-attention decoder
gives the neural operator its discretisation-invariance property. The
Galerkin Transformer \cite{caoChooseTransformer2021} replaces the
softmax attention by a linear attention whose cost scales linearly
with the number of tokens; this is particularly attractive when one
wishes to evaluate the operator on very large meshes. The Transolver
\cite{wuTransolver2024} introduces a small number of physics-aware
latent slots between the input and the queries, so that each slot can
specialise to a different sub-region or physical regime; this is
reported to be advantageous in problems with sharp material contrasts
such as the rod-in-rock setup of Chapter~4. We treat the plain
Transformer of this section as the natural starting point and leave
Galerkin- and Transolver-style refinements as a direction for future
work.

\begin{table}[ht]
\centering
\small
\begin{tabular}{l l l}
\toprule
Architecture & Attention cost & Key idea \\
\midrule
Vanilla Transformer \cite{vaswani2017} & \(\mathcal{O}(N^2 d_m)\) & Softmax dot-product on sequences \\
OFormer \cite{liNeuralOperator2023} & \(\mathcal{O}(N_{\mathrm{in}} N_q d_m)\) & Cross-attention from input set to query set \\
Galerkin Transformer \cite{caoChooseTransformer2021} & \(\mathcal{O}(N d_m^2)\) & Softmax replaced by Galerkin-style linear attention \\
Transolver \cite{wuTransolver2024} & \(\mathcal{O}(N M d_m)\) & Latent physics-aware slots with \(M\!\ll\!N\) \\
\midrule
This work & \(\mathcal{O}(N_{\mathrm{in}} N_q d_m)\) & OFormer baseline, \(N_{\mathrm{in}}\!=\!N_q\!=\!1024\) subsampling \\
\bottomrule
\end{tabular}
\caption{Position of the Transformer-based neural operator used in this
thesis with respect to other set-based neural operators. The notation
\(N\) refers to the total number of tokens, \(N_{\mathrm{in}}\) and
\(N_q\) to the sizes of the input and query sets, \(d_m\) to the
embedding dimension, and \(M\) to the number of latent slots used by
the Transolver.}
\label{tab:transformer-variants}
\end{table}

\subsection{Discretisation invariance and operator-learning view}

Two properties of the construction together make it a neural operator
in the strict sense of the term. First, the model treats the spatial
discretisation as the set of input tokens; an arbitrary subset of the
COMSOL mesh may be used during training without changing the
parameters or the computational graph. Second, the predicted field is
evaluated through a separate set of query coordinates that are not
required to coincide with the inputs; in particular, the model can be
trained on \(1024\) nodes per sample and queried at a much denser
grid at inference. This makes the Transformer-based surrogate a useful
complement to the Fourier neural operator of Section~5.2, which is
intrinsically tied to a regular grid, and to the MeshGraphNet of
Section~5.3, which is intrinsically tied to a particular graph
topology.

\subsection{Summary}

The Transformer-based neural operator of this section is a
discretisation-invariant surrogate for the coupled thermoelastic system
of Chapter~2 trained on the COMSOL ground-truth solutions of
Chapter~4. It tokenises the medium into input points with embedded
material properties and queries it at arbitrary output locations; it
relies on multi-head self-attention in the encoder to capture global
relations within the rock and on multi-head cross-attention in the
decoder to interpolate the predicted thermal and mechanical fields
between mesh nodes; and it is trained with a hybrid loss that mixes
supervised reconstruction, velocity consistency, and a thermal-residual
regulariser. The architecture sits in the OFormer family of neural
operators and is intentionally kept minimal so that its behaviour can
be analysed and compared on an equal footing with the PINN, FNO, and
MeshGraphNet routes evaluated in the remainder of Chapter~5.
